\documentclass[11pt]{article}
\usepackage[T1]{fontenc}
\usepackage[utf8]{inputenc}
\usepackage{lmodern}
\usepackage[margin=1in]{geometry}
\usepackage{amsmath,amssymb,amsthm,mathtools}
\usepackage{enumitem}
\usepackage[hidelinks]{hyperref}
\usepackage{microtype}

\newtheorem{theorem}{Theorem}[section]
\newtheorem{proposition}[theorem]{Proposition}
\newtheorem{lemma}[theorem]{Lemma}
\newtheorem{corollary}[theorem]{Corollary}
\newtheorem{definition}[theorem]{Definition}
\newtheorem{remark}[theorem]{Remark}
\newtheorem{problem}[theorem]{Problem}
\newtheorem{example}[theorem]{Example}

\newcommand{\T}{\mathbb T}
\newcommand{\C}{\mathbb C}
\newcommand{\Q}{\mathbb Q}
\newcommand{\Z}{\mathbb Z}
\newcommand{\R}{\mathbb R}
\DeclareMathOperator{\Zero}{Zero}
\DeclareMathOperator{\Hit}{Hit}
\DeclareMathOperator{\Rad}{rad}
\DeclareMathOperator{\Sat}{sat}
\DeclareMathOperator{\rank}{rank}

\title{Toric Certificates for Effective Skolem Tails and Subspace Orbit Reachability}
\author{Luca Blanchi}
\date{June 2026}

\begin{document}
\maketitle

\begin{abstract}
We develop a certificate framework for proving effective zero-free tails for scalar and vector algebraic linear recurrence sequences. The framework applies to the Skolem Problem, the Simultaneous Skolem Problem, and certified instances of the Subspace Orbit Problem.

After passing to a common Binet representation, splitting root-of-unity degeneracies, saturating residual torsion, and peeling globally cancelled dominant shells, each branch of a vector recurrence is put in the form
\[
\mathbf u_{Mk+r}
=
\Lambda^k\bigl(\mathbf D(k,z(k))+\mathbf E(k)\bigr),
\]
where
\[
z(k)=c\eta^k\in \T^\rho,
\qquad
\mathbf D(T,X)\in L[T][X_1^{\pm1},\ldots,X_\rho^{\pm1}]^m,
\]
and
\[
|\mathbf E(k)|\le C_E(k+1)^s\theta^k,
\qquad
0<\theta<1.
\]
The toric rank, the order of the recurrence, the number of dominant roots, and the root multiplicities are unrestricted.

The main theorem is a soundness theorem. If the dominant vector admits a verified toric lower-bound certificate, then all zeros of the branch lie below a computable bound unless the branch is identically zero. The remaining prefix is checked exactly. We give several verifiable certificate formats: joint semialgebraic certificates, Positivstellensatz and sum-of-squares certificates, Bezout--toric certificates, scalar character-polynomial certificates, and binomial-stratified certificates.

The new automatic layer is binomial-stratified. If the dominant zero locus on the compact toric closure is contained in finitely many effective binomial strata, then a semialgebraic Lojasiewicz inequality gives a lower bound by a power of the distance to those strata. Along the actual toric orbit, Baker--Matveev separation turns the binomial distances into polynomial lower bounds away from exact-hit progressions. Exact hits are computable finite unions of arithmetic progressions; on those progressions one restricts the recurrence and either peels a vanished dominant shell or lowers the toric dimension. In the automatic binomial subclass, the strata are extracted from Laurent elimination ideals and binomial decomposition.

The resulting compiler computes exact zero sets for every recurrence admitting a finite binomial-stratified toric certificate tree, and computes the strata itself in the automatic binomial-stratified class. This isolates a checkable and partially automatic mechanism for proving effective Skolem tails in high-rank, high-order families, including vector examples where scalar coordinatewise certification fails. The geometric boundary of the method is explicit: arbitrary non-binomial toric approximation requires additional lower-bound input.
\end{abstract}

\section*{Declaration of generative AI and AI-assisted technologies in the writing process}
During the preparation of this work, ChatGPT, by OpenAI, was used to assist with mathematical drafting, formalization, review, and editing. This work is shared as a preliminary AI-assisted mathematical note. The mathematical content may have been only partially reviewed and may contain errors; it should not be treated as peer-reviewed or as a fully verified manuscript.

\section{Introduction}

Let \(K\) be a number field. A scalar linear recurrence sequence over \(K\) is a sequence \(u=(u_n)_{n\ge0}\) satisfying
\[
u_{n+d}=a_{d-1}u_{n+d-1}+\cdots+a_0u_n
\]
for some \(a_0,\ldots,a_{d-1}\in K\). The Skolem Problem asks whether there exists \(n\ge0\) such that
\[
u_n=0.
\]
The strong computational form asks to compute the entire zero set
\[
Z(u)=\{n\ge0:u_n=0\}.
\]

The Skolem--Mahler--Lech theorem asserts that \(Z(u)\) is a finite union of arithmetic progressions together with a finite exceptional set. Its known proofs do not provide a general algorithm for computing the finite exceptional set. The computational heart of Skolem is therefore the construction of an effective zero-free tail.

This paper studies such tails through toric certificates. The scalar case is included, but the natural level of generality is vectorial. Let
\[
\mathbf u_n=(u_1(n),\ldots,u_m(n))\in K^m
\]
be a vector of algebraic linear recurrence sequences. The simultaneous zero set is
\[
Z(\mathbf u)=\{n\ge0:\mathbf u_n=0\}.
\]
This formulation covers the Simultaneous Skolem Problem. It also covers affine subspace orbit reachability: if \(A\in K^{d\times d}\), \(x_0\in K^d\), and
\[
V=\{y\in K^d:By=b\},
\]
then
\[
A^n x_0\in V
\quad\Longleftrightarrow\quad
BA^n x_0-b=0.
\]

The central idea is to separate the dominant vector from zero. On each branch, after standard spectral reductions, we obtain
\[
\mathbf u_{Mk+r}
=
\Lambda^k\bigl(\mathbf D(k,z(k))+\mathbf E(k)\bigr),
\]
where \(z(k)\) is a toric orbit on a compact toric coset, \(\mathbf D\) is the normalized dominant vector, and \(\mathbf E\) is exponentially smaller. If one can certify that
\[
\mathbf D(k,z(k))\ne0
\quad\Longrightarrow\quad
|\mathbf D(k,z(k))|
\ge
\exp(-C(\log(k+2))^A)
\]
outside a computable prefix, then the exponentially small error cannot cancel the dominant vector outside that prefix. The rest is an exact finite computation.

The vector viewpoint is essential. A scalar coordinate may contain a factor such as
\[
X+Y-1,
\]
which is outside the automatic one-character Baker class. But a system such as
\[
(X+Y-1,\ X-Y)
\]
has no common zero on \(\T^2\), and the norm of the vector is bounded away from zero by a semialgebraic certificate. Thus the vector method proves effective tails in cases not visible to scalar coordinatewise methods.

The contributions are:
\begin{enumerate}
\item a global vector dominant-shell normal form with torsion saturation;
\item a verifier theorem for toric lower-bound certificates;
\item joint semialgebraic and Positivstellensatz/SOS certificates;
\item Bezout--toric certificates converting scalar lower bounds into vector lower bounds;
\item a scalar Baker--Matveev lemma for expressions \(F(k,\delta^k)\);
\item quotient-character reduction, producing scalar atoms automatically from Laurent factors whose exponent supports become one-dimensional modulo the toric relation lattice;
\item binomial-stratified certificates combining semialgebraic Lojasiewicz bounds, Baker separation from exact binomial hits, and recursive restriction;
\item automatic extraction of binomial strata from dominant elimination ideals in the binomial subclass;
\item a partial Factor--Reduce--Certify compiler and an automatic binomial-stratified certificate-tree compiler;
\item applications to Simultaneous Skolem, Subspace Orbit, affine linear-loop reachability, algebraic target hitting, and eventual cone membership.
\end{enumerate}

The framework isolates certified toric dominant shells: once a dominant lower-bound mechanism is verified, the zero-free tail and the exact finite prefix become effective. Non-binomial toric approximation remains the boundary where additional Diophantine lower bounds would be needed.

\section{Effective input and certificate model}

An input consists of:
\begin{enumerate}
\item a number field \(K\), represented effectively;
\item either a scalar or vector linear recurrence sequence over \(K\), or a linear dynamical system \((A,x_0,B,b)\);
\item a complex embedding \(\sigma:K\hookrightarrow \C\).
\end{enumerate}

All computations take place in finite extensions of \(K\). We use standard exact algorithms in number fields: arithmetic, factorization, construction of splitting fields, root-of-unity testing, exact zero testing, comparison of algebraic absolute values under \(\sigma\), and computation of multiplicative relation lattices \cite{Ge}. These algorithms are not claimed to be efficient in general; they are used as exact symbolic subroutines.

A certificate is extended input. The algorithms in this paper are sound: if the certificate verifies, then the output zero set is correct. They are not complete: failure to certify means ``not certified'', not ``no zero'' and not ``undecidable''.

We separate three layers.
\begin{itemize}
\item The verifier theorem: a valid certificate implies an effective zero-free tail.
\item Automatic certificate generation: for certain structural classes, the certificate is constructed by the compiler.
\item Positive search: for some certificates, such as Bezout--toric identities of bounded degree, one may search positively, but failure of the search has no mathematical consequence.
\end{itemize}

\section{Vector Binet form and global shells}

Let
\[
\mathbf u_n=(u_1(n),\ldots,u_m(n))\in K^m
\]
be a vector of algebraic linear recurrence sequences. Passing to a common splitting field \(L/K\), we may write
\[
\mathbf u_n=\sum_{\lambda\in\mathcal R}\mathbf P_\lambda(n)\lambda^n,
\]
where \(\mathcal R\subset L^\times\) is finite and
\[
\mathbf P_\lambda(T)\in L[T]^m.
\]

Zero characteristic roots contribute only to a finite initial segment. That segment is checked exactly, and the sequence is shifted. Hence all roots in the remaining Binet representation are nonzero.

Fix an embedding
\[
\sigma:L\hookrightarrow\C.
\]
A global vector shell is a set
\[
S_\rho=\{\lambda\in\mathcal R:|\sigma(\lambda)|=\rho\}.
\]
The shell contribution is
\[
\mathbf U_\rho(n)=\sum_{\lambda\in S_\rho}\mathbf P_\lambda(n)\lambda^n.
\]
The shell is vectorially cancelled if
\[
\mathbf U_\rho(n)\equiv0
\]
as a vector-valued exponential-polynomial function. This means every coordinate is identically zero.

After torsion splitting, distinct bases with no root-of-unity quotient give linearly independent exponential-polynomial functions. Therefore vectorial cancellation is decidable exactly.

\section{Torsion splitting and torsion saturation}

If \(\lambda/\mu\) is a root of unity, let \(o(\lambda,\mu)\) be its order. Let
\[
M_0=\operatorname{lcm}\{o(\lambda,\mu):\lambda,\mu\in\mathcal R,\ \lambda/\mu\in\mu_\infty\}.
\]
For each residue class \(r\pmod{M_0}\), define
\[
\mathbf v^{(r)}_k=\mathbf u_{M_0k+r}.
\]
Equal bases in the resulting Binet representation are combined exactly.

A second torsion issue remains. A shell with relative multiplicative rank \(\rho\) modulo torsion need not be an exact toric coset. Its roots have the form
\[
\lambda=\Lambda\,\xi_\lambda\,\gamma^a,
\]
where \(\xi_\lambda\) is a root of unity and
\[
\gamma^a=\gamma_1^{a_1}\cdots\gamma_\rho^{a_\rho},
\qquad
a\in\Z^\rho.
\]
Let
\[
Q=\operatorname{lcm}_{\lambda}\operatorname{ord}(\xi_\lambda).
\]
Passing to subbranches
\[
k=Qh+s
\]
makes
\[
\xi_\lambda^k=\xi_\lambda^s
\]
constant. The polynomial coefficients are transformed by
\[
\mathbf P_\lambda(k)\mapsto \mathbf P_\lambda(Qh+s),
\]
and the bases by
\[
\lambda^k\mapsto \lambda^s(\lambda^Q)^h.
\]
All new data are computable. We call this step torsion saturation.

\section{Vector toric normal form}

After torsion splitting and torsion saturation, peel globally cancelled shells until the first non-cancelled global shell is reached. Let \(S\) be that shell. Choose \(\Lambda\in S\). The roots in \(S\) can be written as
\[
\lambda=\Lambda\gamma^a,
\qquad
a\in A\subset\Z^\rho.
\]
Since all roots in \(S\) have the same \(\sigma\)-absolute value, the toric parameters satisfy
\[
|\sigma(\gamma_j)|=1.
\]

Define
\[
z(k)=(\gamma_1^k,\ldots,\gamma_\rho^k).
\]
Then the contribution of \(S\) has the form
\[
\Lambda^k\mathbf D(k,z(k)),
\]
where
\[
\mathbf D(T,X)\in L[T][X_1^{\pm1},\ldots,X_\rho^{\pm1}]^m.
\]

All lower shells have strictly smaller \(\sigma\)-absolute value. Hence the branch has normal form
\[
\mathbf v_k=\Lambda^k\bigl(\mathbf D(k,z(k))+\mathbf E(k)\bigr),
\]
with
\[
|\mathbf E(k)|\le C_E(k+1)^s\theta^k,
\qquad
0<\theta<1.
\]
Here
\[
|\mathbf w|^2=\sum_{i=1}^m|\sigma(w_i)|^2.
\]

\section{Toric cosets}

After reindexing, it is convenient to write
\[
z(k)=c\eta^k=(c_1\eta_1^k,\ldots,c_\rho\eta_\rho^k),
\]
with
\[
|\sigma(c_j)|=|\sigma(\eta_j)|=1.
\]
Define the exact relation lattice
\[
\mathcal L_\eta
=
\{a=(a_1,\ldots,a_\rho)\in\Z^\rho:\eta^a=1\},
\]
where
\[
\eta^a=\eta_1^{a_1}\cdots\eta_\rho^{a_\rho}.
\]
This lattice is computable.

Let
\[
T_\eta=\{y\in\T^\rho:y^a=1\text{ for all }a\in\mathcal L_\eta\}.
\]
The orbit closure of \(z(k)=c\eta^k\) is the compact toric coset
\[
C=cT_\eta.
\]
Equivalently,
\[
C=\{x\in\T^\rho:x^a=c^a\text{ for all }a\in\mathcal L_\eta\}.
\]

We do not assume \(T_\eta\) is connected. If finite torsion remains in \(\eta\), \(T_\eta\) may have finitely many components. The equations above describe the full compact coset.

For effective real-algebraic verification, write
\[
X_j=x_j+iy_j.
\]
The torus is defined by
\[
x_j^2+y_j^2=1,
\]
and the coset equations are the real and imaginary parts of
\[
X^a=c^a
\]
for a basis of \(\mathcal L_\eta\).

\section{The certified-tail theorem}

\begin{theorem}[Certified-tail soundness]
Let a saturated branch have vector toric normal form
\[
\mathbf v_k=\Lambda^k\bigl(\mathbf D(k,z(k))+\mathbf E(k)\bigr),
\]
where
\[
|\mathbf E(k)|\le C_E(k+1)^s\theta^k,
\qquad
0<\theta<1.
\]
Assume there are computable constants \(A,C,N_0\) such that, for every \(k\ge N_0\),
\[
\mathbf D(k,z(k))\ne0
\quad\Longrightarrow\quad
|\mathbf D(k,z(k))|
\ge
\exp(-C(\log(k+2))^A).
\]
Assume also that exact zeros of
\[
\mathbf D(k,z(k))
\]
are either contained in a computable prefix or occur only on subbranches where the global dominant vector is identically cancelled and hence peeled.

Then the branch has a computable zero-free tail. Consequently the zeros on the branch are computable exactly as a finite set plus any identically zero subprogressions.
\end{theorem}

\begin{proof}
Choose \(B\ge N_0\) such that
\[
C_E(k+1)^s\theta^k
<
\frac12\exp(-C(\log(k+2))^A)
\]
for all \(k\ge B\). Such \(B\) is computable. For \(k\ge B\), if \(\mathbf D(k,z(k))\ne0\), then
\[
|\mathbf D(k,z(k))+\mathbf E(k)|
\ge
|\mathbf D(k,z(k))|-|\mathbf E(k)|
>0.
\]
Hence \(\mathbf v_k\ne0\). The remaining possible zeros are in a computable prefix or in identically zero subbranches produced by peeling. The finite prefix is checked exactly in \(L\).
\end{proof}

\section{Joint semialgebraic certificates}

\begin{definition}[Joint semialgebraic certificate]
Let
\[
\mathbf D(T,X)=(D_1(T,X),\ldots,D_m(T,X)).
\]
A joint semialgebraic certificate, or joint S-certificate, on a toric coset \(C=cT_\eta\) consists of explicit data
\[
B_0\in\mathbb N,\qquad L_0\in\mathbb N,\qquad q\in\Q_{>0}
\]
such that
\[
\sum_{i=1}^m|\sigma(D_i(t,x))|^2
\ge
q(t+1)^{-L_0}
\]
for every
\[
t\ge B_0,
\qquad
x\in C.
\]
\end{definition}

The assertion is a universal formula over a real closed field. Laurent monomials are represented on the torus by
\[
X_j=x_j+iy_j,
\qquad
X_j^{-1}=x_j-iy_j.
\]
The coset \(C\) is represented by polynomial equations
\[
x_j^2+y_j^2=1
\]
and the real and imaginary parts of
\[
X^a=c^a
\]
for a basis of \(\mathcal L_\eta\).

The certificate condition becomes
\[
\forall t,x,y:\ 
\left(t\ge B_0\wedge (x+iy)\in C\right)
\Rightarrow
(t+1)^{L_0}\sum_i|\sigma(D_i(t,x+iy))|^2-q\ge0.
\]
This formula is effectively decidable by quantifier elimination over real closed fields.

\begin{proposition}
A joint S-certificate implies the hypothesis of the certified-tail theorem with a polynomial lower bound:
\[
|\mathbf D(k,z(k))|\ge q^{1/2}(k+1)^{-L_0/2}.
\]
\end{proposition}

\begin{proof}
Since \(z(k)\in C\), substituting \(t=k\) and \(x=z(k)\) into the certificate gives the result.
\end{proof}

\section{Binomial-stratified certificates}

Uniform separation on the whole compact toric coset is a useful certificate, but many natural dominant vectors vanish on structured subtori. The certificate can still be effective when the entire dominant zero locus is forced into finitely many binomial strata.

\begin{definition}[Binomial stratum]
Let \(C\subseteq\T^\rho\) be an effective compact toric coset. A binomial stratum in \(C\) is a subset
\[
H=
\{x\in C:
\chi_{a_1}(x)=\zeta_1,\ldots,\chi_{a_\ell}(x)=\zeta_\ell\},
\]
where
\[
a_j\in\Z^\rho,
\qquad
\chi_{a_j}(X)=X^{a_j},
\]
and each \(\zeta_j\) is algebraic of complex absolute value \(1\). The stratum is proper if \(H\ne C\).
\end{definition}

For such a stratum set
\[
\Delta_H(x)=
\sum_{j=1}^{\ell}|\chi_{a_j}(x)-\zeta_j|^2.
\]
Then
\[
\Delta_H(x)=0
\quad\Longleftrightarrow\quad
x\in H.
\]
For a finite family
\[
\mathcal H=\{H_1,\ldots,H_N\},
\]
define
\[
\Delta_{\mathcal H}(x)=
\prod_{\nu=1}^N\Delta_{H_\nu}(x).
\]
Thus
\[
\Delta_{\mathcal H}(x)=0
\quad\Longleftrightarrow\quad
x\in \bigcup_{\nu=1}^N H_\nu.
\]

\begin{definition}[Binomial stratification certificate]
Let
\[
\mathbf D(T,X)\in L[T][X_1^{\pm1},\ldots,X_\rho^{\pm1}]^m
\]
be a dominant vector on \(C\). A binomial stratification certificate consists of
\[
B\in\mathbb N
\]
and a finite family \(\mathcal H\) of proper effective binomial strata in \(C\) such that
\[
\mathbf D(t,x)=0
\quad\Longrightarrow\quad
x\in\bigcup_{H\in\mathcal H}H
\]
for all real \(t\ge B\) and all \(x\in C\).
\end{definition}

Once the candidate strata are supplied, the implication is a universal formula over a real closed field after writing the torus and the coset in real coordinates. Hence it is effectively checkable by quantifier elimination.

\begin{lemma}[Effective semialgebraic Lojasiewicz inequality]
Let \(K_0\) be a compact effective semialgebraic set, and let
\[
f,g:K_0\to\R_{\ge0}
\]
be effective continuous semialgebraic functions. If
\[
Z(f)\subseteq Z(g),
\]
then there exist effectively searchable constants
\[
N\in\mathbb N,
\qquad
c>0,
\]
with \(c\) algebraic, such that
\[
f(y)\ge c\,g(y)^N
\]
for every \(y\in K_0\).
\end{lemma}

\begin{proof}
The existence statement is the standard semialgebraic Lojasiewicz inequality on compact semialgebraic sets. For effectivity, fix \(N\). The assertion
\[
\exists c>0\ \forall y\in K_0:\ f(y)\ge c\,g(y)^N
\]
is a first-order sentence over real closed fields. Quantifier elimination decides whether such a \(c\) exists and, when it does, gives an effective semialgebraic description of the admissible \(c\)'s. Enumerating \(N=0,1,2,\ldots\) therefore finds a valid exponent and a positive algebraic \(c\).
\end{proof}

\begin{theorem}[Stratified semialgebraic lower bound]
Assume that \(\mathbf D\) admits a binomial stratification certificate \((B,\mathcal H)\) on \(C\). Put
\[
P(t,x)=\sum_{i=1}^m |D_i(t,x)|^2,
\qquad
\Delta(x)=\Delta_{\mathcal H}(x).
\]
Then there exist effectively computable constants
\[
B'\in\mathbb N,
\qquad
q>0,
\qquad
L_0\in\mathbb N,
\qquad
R\in\mathbb N
\]
such that
\[
P(t,x)\ge q(t+1)^{-L_0}\Delta(x)^R
\]
for every \(t\ge B'\) and every \(x\in C\).
\end{theorem}

\begin{proof}
After clearing Laurent monomials, \(P\) is represented by an effective real semialgebraic function on \([B,\infty)\times C\). Choose \(d\) at least the degree in \(t\) of the cleared expression and compactify
\[
u=\frac1{t+1},
\qquad
t=\frac{1-u}{u}.
\]
Set
\[
Q(u,x)=u^dP\left(\frac{1-u}{u},x\right)
\]
on the compact semialgebraic set
\[
K_0=\left[0,\frac1{B+1}\right]\times C.
\]
For \(u>0\), the certificate gives
\[
Q(u,x)=0
\quad\Longrightarrow\quad
\Delta(x)=0.
\]
At \(u=0\), compactification may create additional zeros. Hence
\[
Z(Q)\cap K_0\subseteq Z(u\Delta).
\]
Apply the effective semialgebraic Lojasiewicz inequality to
\[
f=Q,
\qquad
g=u\Delta.
\]
There are \(N\) and \(c>0\) such that
\[
Q(u,x)\ge c\,u^N\Delta(x)^N.
\]
Since \(Q=u^dP\), this gives
\[
P(t,x)\ge c\,u^{N-d}\Delta(x)^N.
\]
If \(N\ge d\), take \(L_0=N-d\). If \(N<d\), weaken the estimate on \(0<u\le1\) to \(P(t,x)\ge c\Delta(x)^N\). This gives the stated form with effective constants.
\end{proof}

\section{SOS and Positivstellensatz certificates}

Quantifier elimination proves decidability, but one may also use algebraic certificates that are easier to verify.

Let \(C\) be represented by equalities
\[
f_1=\cdots=f_r=0
\]
and inequalities
\[
g_1\ge0,\ldots,g_s\ge0.
\]
These include the torus equations, the coset equations, and the tail condition \(t-B_0\ge0\). A Positivstellensatz/SOS certificate for the joint lower bound consists of an identity
\[
(t+1)^L\sum_i |D_i(t,X)|^2-q
=
\sigma_0+\sum_j\sigma_j g_j+\sum_\ell h_\ell f_\ell,
\]
where the \(\sigma_j\) are sums of squares and the \(h_\ell\) are arbitrary polynomials, all with algebraic coefficients.

Such an identity is checked by exact polynomial arithmetic and by verifying the supplied sum-of-squares decompositions. When present, it gives a joint S-certificate.

This SOS formulation is not required for decidability, but it gives a more explicit certificate format than quantifier elimination.

\section{Automatic joint certificates from leading vectors}

Write
\[
\mathbf D(T,X)=T^d\mathbf D_d(X)+T^{d-1}\mathbf D_{d-1}(X)+\cdots+\mathbf D_0(X),
\]
with \(\mathbf D_d\ne0\).

\begin{proposition}[Leading-vector criterion]
If
\[
\mathbf D_d(x)\ne0
\qquad\text{for every }x\in C,
\]
then a joint S-certificate exists and can be effectively found.
\end{proposition}

\begin{proof}
The function
\[
x\mapsto |\mathbf D_d(x)|^2
\]
is continuous and positive on the compact coset \(C\). Its positivity can be verified by quantifier elimination. Once verified, one can find a rational \(q_0>0\) such that
\[
|\mathbf D_d(x)|^2\ge q_0
\]
for all \(x\in C\).

For \(j<d\), compute an upper bound \(M_j\) for \(|\mathbf D_j(x)|\) on \(C\). Since monomials have modulus \(1\) on the torus, summing absolute values of algebraic coefficients gives an effective upper bound.

For \(t\) large enough,
\[
\left|\sum_{j<d}t^j\mathbf D_j(x)\right|
\le
\frac12q_0^{1/2}t^d
\]
uniformly on \(C\). Hence
\[
|\mathbf D(t,x)|\ge \frac12q_0^{1/2}t^d.
\]
This is a joint S-certificate.
\end{proof}

\section{Scalar character-polynomial atoms}

A scalar atom is an expression
\[
F(T,\chi(X)),
\]
where
\[
F(T,Y)\in L[T,Y]\setminus\{0\}
\]
and
\[
\chi(X)=X_1^{a_1}\cdots X_\rho^{a_\rho}
\]
is a toric character. Along the branch
\[
z(k)=c\eta^k,
\]
we have
\[
\chi(z(k))=\chi(c)\chi(\eta)^k.
\]
Absorbing the constant \(\chi(c)\) into \(F\), one reduces to
\[
F(k,\delta^k),
\qquad
\delta=\chi(\eta).
\]

\section{The Baker--Matveev lemma}

We use the following standard consequence of effective lower bounds for linear forms in logarithms.

\begin{lemma}[Scalar character lower bound]
Let
\[
F(T,Y)\in L[T,Y]\setminus\{0\},
\]
let
\[
\delta\in L^\times,
\qquad
|\sigma(\delta)|=1,
\]
and suppose \(\delta\) is not a root of unity. Then there exist computable constants
\[
C>0,
\qquad
N_0\in\mathbb N
\]
such that, for every \(k\ge N_0\),
\[
F(k,\delta^k)\ne0
\quad\Longrightarrow\quad
|\sigma(F(k,\delta^k))|
\ge
\exp(-C(\log(k+2))^2).
\]
Moreover, exact zeros of \(F(k,\delta^k)\) are contained in a computable finite prefix.
\end{lemma}

\begin{proof}
Write
\[
F(T,Y)=\sum_{j=0}^r f_j(T)Y^j.
\]
For each \(k\), let
\[
F_k(Y)=F(k,Y).
\]
The set of \(k\) such that \(F_k\equiv0\) is finite and computable unless all \(f_j\) vanish identically, which is excluded. The case in which \(F_k\) is constant nonzero is harmless and gives a Liouville lower bound.

Assume \(F_k\) is nonconstant. Factor
\[
F_k(Y)=a_k\prod_\nu(Y-\eta_{\nu,k})
\]
over \(\overline{\Q}\). The degree of \(F_k\) is at most \(r\), though it may drop for special \(k\). This only decreases the number of factors.

The coefficients \(f_j(k)\) have logarithmic height \(O(\log(k+2))\). Standard height estimates for roots of a polynomial give
\[
h(\eta_{\nu,k})\le C_1\log(k+2)
\]
with \(C_1\) effective. Liouville's inequality gives
\[
|\sigma(a_k)|\ge\exp(-C_2\log(k+2))
\]
whenever \(a_k\ne0\).

Let \(\eta=\eta_{\nu,k}\). If \(\eta=0\), then
\[
|\sigma(\delta^k-\eta)|=1.
\]
Assume \(\eta\ne0\) and \(\delta^k\ne\eta\). Apply an effective lower bound for linear forms in logarithms to
\[
\eta^{-1}\delta^k-1.
\]
The degree of the field \(\Q(\delta,\eta)\) is uniformly bounded in terms of \(L\) and \(r\). The height of \(\eta\) is \(O(\log k)\), and the exponent of \(\delta\) is \(k\). A standard Baker--Wuestholz or Matveev estimate yields
\[
|\sigma(\eta^{-1}\delta^k)-1|
\ge
\exp(-C_3(\log(k+2))^2).
\]
Liouville gives
\[
|\sigma(\eta)|\ge\exp(-C_4\log(k+2)).
\]
Therefore
\[
|\sigma(\delta^k-\eta)|
\ge
\exp(-C_5(\log(k+2))^2).
\]

If \(F(k,\delta^k)\ne0\), no factor vanishes. Multiplying the lower bounds for the leading coefficient and all nonzero factors gives
\[
|\sigma(F(k,\delta^k))|
\ge
\exp(-C(\log(k+2))^2).
\]

Now suppose
\[
F(k,\delta^k)=0.
\]
Then \(\delta^k\) is a root of \(F_k\), so
\[
h(\delta^k)\le C_1\log(k+2).
\]
Since \(\delta\) is not a root of unity, \(h(\delta)>0\), and
\[
h(\delta^k)=k h(\delta).
\]
Thus
\[
k h(\delta)\le C_1\log(k+2),
\]
which fails beyond a computable bound.
\end{proof}

\section{Baker separation from binomial strata}

The preceding scalar lemma controls one-character atoms. Binomial strata reduce to finitely many such one-character conditions along the toric orbit.

\begin{lemma}[One-character separation]
Let
\[
\alpha,\zeta\in\overline{\Q}^{\times},
\qquad
|\sigma(\alpha)|=|\sigma(\zeta)|=1.
\]
Then the solution set
\[
S(\alpha,\zeta)=\{k\ge0:\alpha^k=\zeta\}
\]
is effectively computable as a finite union of arithmetic progressions and a finite set. Moreover, there exist effective constants
\[
c>0,
\qquad
A\in\mathbb N,
\]
such that
\[
|\sigma(\alpha^k-\zeta)|\ge c(k+1)^{-A}
\]
for every \(k\notin S(\alpha,\zeta)\).
\end{lemma}

\begin{proof}
If \(\alpha\) is a root of unity, the sequence \(\alpha^k\) is periodic. Thus \(S(\alpha,\zeta)\) is empty or a finite union of residue classes modulo the order of \(\alpha\), and the nonzero distances outside \(S(\alpha,\zeta)\) have a positive periodic minimum.

Assume \(\alpha\) is not a root of unity. Then \(\alpha^k=\zeta\) has at most one solution, since two distinct solutions would make \(\alpha\) torsion. The possible exact solution is decidable by the standard algorithms for multiplicative relations among algebraic numbers. For non-solutions, apply an effective lower bound for linear forms in logarithms to
\[
\alpha^k\zeta^{-1}-1.
\]
Since \(\alpha\) and \(\zeta\) are fixed, Baker--Wuestholz or Matveev estimates give a lower bound of the form
\[
|\sigma(\alpha^k\zeta^{-1}-1)|\ge c_1(k+1)^{-A}.
\]
Multiplying by the fixed number \(|\sigma(\zeta)|=1\) gives the stated estimate.
\end{proof}

\begin{lemma}[Separation from one binomial stratum]
Let \(z(k)=c\eta^k\in C\), and let
\[
H=
\{x\in C:\chi_{a_1}(x)=\zeta_1,\ldots,\chi_{a_\ell}(x)=\zeta_\ell\}
\]
be an effective binomial stratum. Set
\[
S_H=\{k\ge0:z(k)\in H\}.
\]
Then \(S_H\) is effectively computable as a finite union of arithmetic progressions and a finite set. Moreover, there exist effective constants
\[
c_H>0,
\qquad
A_H\in\mathbb N,
\]
such that
\[
\Delta_H(z(k))\ge c_H(k+1)^{-A_H}
\]
for every \(k\notin S_H\).
\end{lemma}

\begin{proof}
The condition \(z(k)\in H\) is the finite system
\[
\chi_{a_j}(\eta)^k
=
\zeta_j\,\chi_{a_j}(c)^{-1},
\qquad
j=1,\ldots,\ell.
\]
Each equation has an effectively computable solution set by the one-character lemma, and finite intersections of finite unions of progressions and finite sets are effectively computable.

If \(k\notin S_H\), then at least one equation fails. For each \(j\), the one-character lemma gives a polynomial lower bound away from the exact solution set of the \(j\)-th equation. Taking the minimum over the finitely many failed possibilities gives
\[
\Delta_H(z(k))
=
\sum_j|\chi_{a_j}(z(k))-\zeta_j|^2
\ge c_H(k+1)^{-A_H}.
\]
\end{proof}

\begin{lemma}[Separation from a finite family of binomial strata]
Let \(\mathcal H=\{H_1,\ldots,H_N\}\) be a finite family of effective binomial strata and set
\[
S_{\mathcal H}=\bigcup_{\nu=1}^N S_{H_\nu}.
\]
Then \(S_{\mathcal H}\) is effectively computable as a finite union of arithmetic progressions and a finite set. Moreover, there exist effective constants
\[
c>0,
\qquad
A\in\mathbb N,
\]
such that
\[
\Delta_{\mathcal H}(z(k))\ge c(k+1)^{-A}
\]
for every \(k\notin S_{\mathcal H}\).
\end{lemma}

\begin{proof}
The exact-hit set is a finite union of effectively computable sets. If \(k\notin S_{\mathcal H}\), then \(z(k)\) lies in none of the strata. Hence every \(\Delta_{H_\nu}(z(k))\) is bounded below by a negative power of \(k+1\). Multiplying the finitely many inequalities gives the stated lower bound for
\[
\Delta_{\mathcal H}(z(k))
=
\prod_{\nu=1}^N\Delta_{H_\nu}(z(k)).
\]
\end{proof}

\begin{theorem}[Stratified zero-free tail outside exact hits]
Let a saturated branch have normal form
\[
\mathbf v_k=\Lambda^k(\mathbf D(k,z(k))+\mathbf E(k)),
\qquad
z(k)=c\eta^k,
\]
with
\[
|\mathbf E(k)|\le C_E(k+1)^s\theta^k,
\qquad
0<\theta<1.
\]
Assume that \(\mathbf D\) admits a binomial stratification certificate \(\mathcal H\) on \(C\), and let
\[
S_{\mathcal H}=\{k\ge0:z(k)\in\bigcup_{H\in\mathcal H}H\}.
\]
Then there is an effectively computable \(B\) such that
\[
\mathbf v_k\ne0
\]
for every \(k\ge B\) with \(k\notin S_{\mathcal H}\).
\end{theorem}

\begin{proof}
The stratified semialgebraic lower bound gives
\[
\sum_i|D_i(k,z(k))|^2
\ge
q(k+1)^{-L_0}\Delta_{\mathcal H}(z(k))^R.
\]
For \(k\notin S_{\mathcal H}\), Baker separation from the finite family of strata gives
\[
\Delta_{\mathcal H}(z(k))\ge c(k+1)^{-A}.
\]
Therefore
\[
|\mathbf D(k,z(k))|
\ge
c'(k+1)^{-A'}
\]
for effective constants \(c'>0\) and \(A'\). The exponentially small error is eventually less than half of this lower bound, and the conclusion follows from the same comparison used in the certified-tail theorem.
\end{proof}

\section{Torsion scalar atoms}

If \(\delta\) is a root of unity of order \(Q\), pass to subbranches
\[
k=Qh+s.
\]
Then
\[
F(k,\delta^k)=F(Qh+s,\delta^s)=G_s(h),
\]
where \(G_s\in L[h]\).

If \(G_s\not\equiv0\), its integer zeros are finite and computable, and Liouville gives a polynomial lower bound for nonzero values.

If \(G_s\equiv0\), the scalar atom vanishes identically on that subbranch. In a scalar product certificate this may force the dominant scalar to vanish. In a vector or Bezout certificate it means only that the particular scalar test has failed. The algorithm either peels the branch if the full dominant vector is identically zero, or requires a new certificate on that subbranch.

Termination follows because every peeling removes one global shell and all splitting operations are finite.

\section{Bezout--toric certificates}

Let
\[
\mathbf D=(D_1,\ldots,D_m).
\]
Let
\[
a^{(1)},\ldots,a^{(\ell)}
\]
be a basis of \(\mathcal L_\eta\), and define the coset binomials
\[
G_j(X)=X^{a^{(j)}}-c^{a^{(j)}}.
\]
Then
\[
G_j(z(k))=0.
\]

\begin{definition}[Bezout--toric certificate]
A Bezout--toric certificate consists of Laurent polynomials
\[
A_i(T,X),\quad B_j(T,X),\quad M(T,X)
\]
such that
\[
M(T,X)=\sum_{i=1}^m A_i(T,X)D_i(T,X)+\sum_{j=1}^\ell B_j(T,X)G_j(X),
\]
together with a scalar certificate for \(M(k,z(k))\) giving
\[
|M(k,z(k))|\ge \exp(-C(\log(k+2))^A)
\]
outside a computable prefix.
\end{definition}

The identity is verified exactly in the Laurent polynomial ring.

\begin{lemma}
There exist computable constants \(C_A,R\) such that
\[
\left(\sum_i|A_i(k,z(k))|^2\right)^{1/2}
\le C_A(k+1)^R.
\]
\end{lemma}

\begin{proof}
Write
\[
A_i(T,X)=\sum_\nu a_{i,\nu}(T)X^{b_{i,\nu}}.
\]
On the torus,
\[
|X^{b_{i,\nu}}|=1.
\]
Each polynomial \(a_{i,\nu}(k)\) grows at most polynomially in \(k\). Summing finitely many terms gives the stated bound.
\end{proof}

\begin{proposition}
A Bezout--toric certificate implies the lower-bound hypothesis of the certified-tail theorem.
\end{proposition}

\begin{proof}
Along the orbit,
\[
M(k,z(k))=\sum_i A_i(k,z(k))D_i(k,z(k)).
\]
By Cauchy--Schwarz,
\[
|M(k,z(k))|
\le
\left(\sum_i|A_i(k,z(k))|^2\right)^{1/2}|\mathbf D(k,z(k))|.
\]
Using the preceding lemma,
\[
|M(k,z(k))|\le C_A(k+1)^R|\mathbf D(k,z(k))|.
\]
A subexponential lower bound for \(M\) therefore gives a subexponential lower bound for \(|\mathbf D|\).
\end{proof}

One may search for Bezout--toric certificates by fixing degree bounds for the unknown \(A_i,B_j,M\), imposing the identity as a finite linear system over \(L\), and checking whether \(M\) admits a scalar S/B/T certificate. If the search succeeds, the certificate is valid. If it fails at a chosen bound, no conclusion follows. Thus this is a positive semidecision procedure, not a complete algorithm.

\section{Quotient-character reduction}

Let
\[
H(T,X)=\sum_{a\in A}h_a(T)X^a
\]
be a Laurent factor. Consider the image of \(A\) in
\[
\Z^\rho/\mathcal L_\eta.
\]
We say that \(H\) has quotient affine rank at most one if the differences
\[
a-a_0,
\qquad
a\in A,
\]
span a vector space of dimension at most one in
\[
(\Z^\rho/\mathcal L_\eta)\otimes_\Z\Q.
\]

\begin{proposition}
If \(H\) has quotient affine rank at most one, then after finite torsion splitting,
\[
H(k,z(k))=\alpha_0^k F(k,\delta^k)
\]
for some algebraic \(\alpha_0,\delta\) of \(\sigma\)-absolute value \(1\) and some
\[
F(T,Y)\in L[T,Y].
\]
\end{proposition}

\begin{proof}
Choose \(a_0\in A\). Since the quotient affine rank is at most one, there exists \(v\in\Z^\rho\) such that every \(a-a_0\) maps to a rational multiple of \(v\) in the quotient. Clearing denominators and splitting finite torsion in the quotient, we may assume that for each \(a\in A\),
\[
a=a_0+r_av+\ell_a
\]
with \(r_a\in\Z\) and \(\ell_a\in\mathcal L_\eta\).

On the coset \(z(k)=c\eta^k\),
\[
z(k)^a
=
c^a\eta^{ak}
=
c^a\eta^{a_0k}(\eta^v)^{r_ak},
\]
because \(\eta^{\ell_a}=1\). Thus
\[
H(k,z(k))
=
\eta^{a_0k}\sum_{a\in A}h_a(k)c^a(\eta^v)^{r_ak}.
\]
After multiplying by a power of \(Y\) if necessary, the sum becomes \(F(k,\delta^k)\) with
\[
\delta=\eta^v.
\]
The factor \(\eta^{a_0k}\) is \(\alpha_0^k\) and has modulus \(1\).
\end{proof}

\begin{corollary}
A quotient-affine-rank-one factor is automatically converted into a scalar Baker or torsion atom.
\end{corollary}

\section{Automatic binomial extraction}

The binomial-stratified theorem is a verifier: it proves a tail once the strata are known. In an important subclass the strata can be extracted from the dominant equations.

\subsection{Dominant elimination ideals}

Let \(C\subseteq\T^\rho\) be the compact toric coset attached to a branch, and let
\[
I_C\subseteq L[X_1^{\pm1},\ldots,X_\rho^{\pm1}]
\]
be its Laurent ideal, generated by the binomials
\[
X^a-c^a,
\qquad
a\in\mathcal L_\eta.
\]
For
\[
\mathbf D(T,X)=(D_1,\ldots,D_m)
\]
define
\[
I_D=
I_C+\langle D_1,\ldots,D_m\rangle
\subseteq
L[T,X_1^{\pm1},\ldots,X_\rho^{\pm1}]
\]
and
\[
J_D=
\Rad\left(I_D\cap L[X_1^{\pm1},\ldots,X_\rho^{\pm1}]\right).
\]
The variety \(V(J_D)\cap C\) contains every point \(x\in C\) for which there exists a complex value \(T=t\) with
\[
\mathbf D(t,x)=0.
\]
It is therefore an algebraic over-approximation of the eventual real projection
\[
\{x\in C:\exists t\gg0,\ \mathbf D(t,x)=0\}.
\]
The over-approximation is useful: if it is already covered by proper binomial strata, then it gives a valid certificate without guessing the strata.

\begin{definition}[Automatically binomial branch]
A dominant branch is automatically binomial if the computed set
\[
V(J_D)\cap C
\]
is effectively verified to be contained in a finite union of proper binomial strata
\[
H_1\cup\cdots\cup H_N
\]
inside \(C\). In the strict binomial case, these strata are obtained from binomial radical and decomposition algorithms \cite{EisenbudSturmfels}, followed by lattice basis reduction and Smith normal form.
\end{definition}

\begin{theorem}[Automatic binomial certificate extraction]
Let a nonzero dominant branch be automatically binomial. Then one can effectively compute a finite family of proper binomial strata
\[
\mathcal H=\{H_1,\ldots,H_N\}
\]
such that
\[
\mathbf D(t,x)=0
\quad\Longrightarrow\quad
x\in\bigcup_{\nu=1}^N H_\nu
\]
for every \(t\) and every \(x\in C\). Hence the branch admits a binomial stratification certificate with \(B=0\).
\end{theorem}

\begin{proof}
Compute \(I_C\), form \(I_D\), eliminate \(T\), and take the radical \(J_D\). By the automatic binomial hypothesis, \(V(J_D)\cap C\) is covered by effectively listed proper binomial strata. If \(\mathbf D(t,x)=0\) for some \(x\in C\), then \((t,x)\in V(I_D)\). Every Laurent polynomial in \(I_D\cap L[X_1^{\pm1},\ldots,X_\rho^{\pm1}]\) therefore vanishes at \(x\), and so \(x\in V(J_D)\cap C\). The verified binomial cover gives the desired implication.
\end{proof}

\section{Binomial-stratified certificate trees}

Exact hits of binomial strata are not exceptional failures. They are arithmetic progressions on which the recurrence can be restricted and renormalized.

\begin{definition}[Binomial-stratified toric certificate tree]
A binomial-stratified toric certificate tree for a vector recurrence is a finite rooted tree whose nodes carry restricted branches
\[
\mathbf v_k=\mathbf u_{ak+b}
\]
with saturated dominant normal forms
\[
\mathbf v_k=\Lambda^k(\mathbf D(k,z(k))+\mathbf E(k)).
\]

A terminal node is one of the following:
\begin{enumerate}
\item an identically zero branch;
\item a branch with a joint S-certificate, an SOS/Positivstellensatz certificate, a Bezout--toric certificate, or scalar atoms proving a dominant lower bound;
\item a closure-avoiding branch for which the leading-vector or semialgebraic compactness test gives a lower bound.
\end{enumerate}

A nonterminal node carries a binomial stratification certificate \(\mathcal H\). Its children are the infinite arithmetic progressions contained in
\[
S_{\mathcal H}=
\{k\ge0:z(k)\in\bigcup_{H\in\mathcal H}H\}.
\]
On each child, the recurrence is restricted to that progression and the dominant normal form is recomputed. Along every edge, either a dominant spectral shell is peeled or the dimension of the compact toric orbit closure strictly decreases.
\end{definition}

\begin{definition}[Automatic binomial-stratified tree]
A binomial-stratified toric certificate tree is automatic if every nonterminal node obtains its binomial stratification certificate from automatic binomial extraction.
\end{definition}

\begin{theorem}[Exact zero set from a certificate tree]
If a vector recurrence admits a finite binomial-stratified toric certificate tree, then its zero set is computable exactly as a finite union of arithmetic progressions and a finite set. If the tree is automatic, the binomial strata at all nonterminal nodes are computed from the dominant branch equations.
\end{theorem}

\begin{proof}
Proceed by induction on the finite tree. An identically zero terminal node contributes its whole arithmetic progression. Any other terminal node has a verified dominant lower bound, hence a computable zero-free tail by the certified-tail theorem; its finite prefix is checked exactly.

At a nonterminal node, the stratified zero-free tail theorem gives a computable bound beyond which no zeros occur outside \(S_{\mathcal H}\). The finitely many remaining indices outside \(S_{\mathcal H}\), together with the finite exceptional part of \(S_{\mathcal H}\), are checked exactly. The infinite progressions in \(S_{\mathcal H}\) are precisely the children, whose zero sets are computed by induction. Taking the union over all nodes gives the exact zero set.
\end{proof}

\section{The Factor--Reduce--Certify compiler}

The partial compiler operates branchwise.
\begin{enumerate}
\item Compute the vector toric normal form.
\item Compute the toric coset \(C=cT_\eta\).
\item Attempt the leading-vector criterion.
\item If it fails, attempt to construct a Bezout--toric certificate by bounded-degree linear algebra.
\item For scalar factors or scalar tests, factor in the Laurent polynomial ring.
\item For each irreducible factor, compute its quotient affine rank modulo \(\mathcal L_\eta\).
\item If the rank is at most one, reduce it to a scalar atom.
\item Verify remaining S-certificates by quantifier elimination, SOS identities, or Positivstellensatz identities.
\item If the branch is not terminally certified, attempt automatic binomial extraction from \(J_D\).
\item If binomial extraction succeeds, compute the exact-hit progressions, restrict to the infinite progressions, and recompute the dominant normal form on each child.
\item Accept the recursive step only when it peels a dominant shell or strictly lowers the toric dimension.
\item If every terminal branch is certified and every recursive branch belongs to a finite certificate tree, compute the zero-free tails and exact prefixes.
\item If some branch is not certified, return ``not certified''.
\end{enumerate}

\begin{theorem}[Compiler soundness]
If Factor--Reduce--Certify succeeds on all branches of a vector recurrence \(\mathbf u\), then it computes
\[
Z(\mathbf u)=\{n\ge0:\mathbf u_n=0\}
\]
exactly as a finite union of arithmetic progressions and a finite set.
\end{theorem}

\begin{proof}
Every terminal branch has one of the terminal certificates listed above, so it satisfies the certified-tail theorem. Identically zero branches give progressions. Nonzero terminal branches have zero-free tails and exact finite prefix checks.

Every accepted nonterminal step is justified by a binomial stratification certificate. The stratified zero-free tail theorem removes all sufficiently large indices outside the exact-hit progressions. The children are precisely the restrictions to those infinite progressions. Since a successful run produces a finite certificate tree, the exact zero set is computed by the theorem on certificate trees.
\end{proof}

\section{Certified and automatic binomial-stratified classes}

\begin{definition}
A vector recurrence belongs to the certified binomial-stratified toric class if its dominant branch decomposition admits a finite binomial-stratified toric certificate tree.

It belongs to the automatic binomial-stratified toric class if it admits such a tree and every nonterminal stratum is produced by automatic binomial extraction.
\end{definition}

\begin{corollary}[Certified binomial-stratified Skolem]
For every algebraic vector recurrence \(\mathbf u\) in the certified binomial-stratified toric class, the set
\[
Z(\mathbf u)=\{n\ge0:\mathbf u_n=0\}
\]
is computable exactly.
\end{corollary}

\begin{proof}
Apply the certificate-tree theorem to every residue branch in the dominant toric normal form and take the finite union of the branch zero sets.
\end{proof}

\begin{corollary}[Automatic binomial-stratified Skolem]
For every algebraic vector recurrence \(\mathbf u\) in the automatic binomial-stratified toric class, the zero set \(Z(\mathbf u)\) is computable exactly. In this class the binomial strata required at nonterminal nodes are computed from the dominant elimination ideals.
\end{corollary}

\begin{remark}
The restriction is geometric rather than dimensional. There is no a priori bound here on recurrence order, root multiplicity, number of dominant roots, toric rank, vector dimension, ambient matrix dimension, or target dimension. The required structure is that the dominant toric zero loci are either terminally separated from zero or recursively controlled by binomial strata.
\end{remark}

\section{Applications}

\begin{theorem}[Simultaneous Skolem]
Let
\[
u^{(1)},\ldots,u^{(m)}
\]
be algebraic linear recurrence sequences and set
\[
\mathbf u_n=(u^{(1)}_n,\ldots,u^{(m)}_n).
\]
If Factor--Reduce--Certify succeeds on \(\mathbf u\), then
\[
\{n\ge0:u^{(1)}_n=\cdots=u^{(m)}_n=0\}
\]
is computable exactly.
\end{theorem}

This is stronger than separately solving the scalar Skolem instances and intersecting their zero sets, because a joint vector certificate may exist even when a scalar coordinate is outside the automatic scalar certificate class.

\begin{theorem}[Subspace Orbit and affine linear loops]
Let
\[
A\in K^{d\times d},
\qquad
x_0\in K^d,
\]
and let
\[
V=\{y\in K^d:By=b\}
\]
be an affine subspace. Define
\[
\mathbf u_n=BA^nx_0-b.
\]
If Factor--Reduce--Certify succeeds on \(\mathbf u_n=BA^nx_0-b\), then the hitting-time set
\[
\{n\ge0:A^nx_0\in V\}
\]
is computable exactly.
\end{theorem}

\begin{proof}
We have
\[
A^nx_0\in V
\quad\Longleftrightarrow\quad
BA^nx_0-b=0.
\]
Apply compiler soundness.
\end{proof}

\begin{corollary}[Binomial-stratified Subspace Orbit]
With notation as above, suppose
\[
\mathbf u_n=BA^nx_0-b
\]
belongs to the certified binomial-stratified toric class. Then
\[
\Hit(A,x_0,V)=\{n\ge0:A^nx_0\in V\}
\]
is computable exactly. If \(\mathbf u_n\) belongs to the automatic binomial-stratified toric class, the required binomial strata are computed from the dominant branch equations.
\end{corollary}

\begin{proof}
The sequence \(\mathbf u_n\) is a vector recurrence, and
\[
A^nx_0\in V
\quad\Longleftrightarrow\quad
\mathbf u_n=0.
\]
Apply the certified or automatic binomial-stratified Skolem corollary.
\end{proof}

This gives a certified reachability theorem for affine linear loops
\[
x:=Ax
\]
with target condition \(Bx=b\).

\begin{corollary}[Algebraic target hitting]
Let
\[
W=\{y:P_1(y)=\cdots=P_m(y)=0\}
\]
be an affine algebraic set over \(K\). Since coordinates of \(A^nx_0\) are linear recurrence sequences and such sequences are closed under addition and pointwise multiplication, each sequence
\[
P_i(A^nx_0)
\]
is a linear recurrence sequence. If Factor--Reduce--Certify succeeds on
\[
\bigl(P_1(A^nx_0),\ldots,P_m(A^nx_0)\bigr),
\]
then
\[
\{n\ge0:A^nx_0\in W\}
\]
is computable exactly.
\end{corollary}

\section{Eventual cone membership and ultimate positivity}

Let \(C\subseteq\R^m\) be a semialgebraic cone or a closed semialgebraic set. Suppose a real vector recurrence has normal form
\[
\mathbf u_k=\Lambda^k(\mathbf D(k,z(k))+\mathbf E(k))
\]
with \(\Lambda>0\), and suppose a certificate proves that \(\mathbf D(k,z(k))\) has distance at least
\[
c(k+1)^{-L}
\]
from the boundary of \(C\) for \(k\ge B\), with a fixed side of the boundary. Then \(\mathbf u_k\in C\) eventually if and only if the dominant vector lies on that side. The finite prefix is checked exactly.

For \(m=1\) and \(C=[0,\infty)\), this gives certified ultimate positivity.

\section{A high-rank vector family}

Let \(\rho\ge2\), and let
\[
\gamma_1,\ldots,\gamma_\rho\in L^\times
\]
satisfy
\[
|\sigma(\gamma_i)|=1
\]
and be multiplicatively independent modulo torsion.

Define
\[
\mathbf D_\rho(X_1,\ldots,X_\rho)
=
\left(
X_2-X_1,\,
X_3-X_1,\,
\ldots,\,
X_\rho-X_1,\,
X_1+\cdots+X_\rho-1
\right).
\]

\begin{proposition}
The vector \(\mathbf D_\rho\) has no zero on \(\T^\rho\).
\end{proposition}

\begin{proof}
If \(\mathbf D_\rho(X)=0\), then
\[
X_2=X_1,\quad \ldots,\quad X_\rho=X_1.
\]
Thus
\[
X_1=\cdots=X_\rho.
\]
The last component gives
\[
\rho X_1=1.
\]
But \(|X_1|=1\), whereas \(|1/\rho|<1\) for \(\rho\ge2\). This is impossible.
\end{proof}

Therefore
\[
Q_\rho(X)=|\mathbf D_\rho(X)|^2
\]
has a strictly positive minimum on \(\T^\rho\), and a joint S-certificate can be verified.

Let
\[
\mathbf u_n
=
\Lambda^n\left(\mathbf D_\rho(\gamma_1^n,\ldots,\gamma_\rho^n)+\mathbf E(n)\right),
\]
where
\[
|\mathbf E(n)|\le C\theta^n,
\qquad
0<\theta<1.
\]
Then Factor--Reduce--Certify succeeds and computes the simultaneous zero set.

The scalar coordinate
\[
X_1+\cdots+X_\rho-1
\]
has exponent support of affine rank \(\rho\), so it is not an automatic quotient-character scalar factor. The vector system is nevertheless certified by joint separation. This demonstrates the genuine advantage of the vector method.

\section{Realization as Subspace Orbit}

Let
\[
A=\operatorname{diag}(\gamma_1,\ldots,\gamma_\rho),
\qquad
x_0=(1,\ldots,1).
\]
Then
\[
A^nx_0=(\gamma_1^n,\ldots,\gamma_\rho^n).
\]
Let \(V_\rho\subseteq K^\rho\) be the affine subspace defined by
\[
x_2-x_1=0,\quad
x_3-x_1=0,\quad\ldots,\quad x_\rho-x_1=0,
\]
and
\[
x_1+\cdots+x_\rho=1.
\]
Then
\[
A^nx_0\in V_\rho
\quad\Longleftrightarrow\quad
\mathbf D_\rho(\gamma_1^n,\ldots,\gamma_\rho^n)=0.
\]
By the preceding proposition, the pure dominant model has no hitting time. With exponentially smaller subdominant perturbations, the subspace-orbit certificate theorem gives a computable zero-free tail and exact prefix check.

This gives a high-rank family of certified Subspace Orbit instances in which one scalar equation contains high-rank toric support, but the vector system is jointly separated.

\section{Relation to previous work}

The Skolem Problem is open in general. The Skolem--Mahler--Lech theorem gives the eventual structure of the zero set but is not a general algorithm for computing the finite exceptional set.

Low-order and few-dominant-root results decide important special cases. The present framework is different: it is not based on low order or few dominant roots, but on a verifiable toric separation certificate.

Modern algorithms for broader recurrence classes often produce correct certificates when they terminate, with termination proved under strong Diophantine conjectures. The present method is unconditional inside the certified class.

The Evertse--Schlickewei--Schmidt theory gives explicit upper bounds for the number of nondegenerate solutions of linear equations in variables from a finite-rank multiplicative group \cite{ESS}. Such results are close in spirit, but a bound on the number of solutions does not by itself give a computable bound on the positions of the indices \(n\). The present method produces an explicit zero-free tail once a certificate is verified.

Bacik and Varonka study Subspace Orbit and Simultaneous Skolem from a dimension-theoretic perspective \cite{BacikVaronka}. Their results are orthogonal to the present work. They prove decidability in certain target-dimension regimes and hardness in others. The present framework imposes no a priori target-dimension threshold; it requires a toric separation certificate for the dominant vector.

The vector feature is essential. The family above shows that scalar coordinatewise certification can fail while joint vector certification succeeds.

\section{Boundary of the certified class}

The certified class is governed by the available lower-bound mechanisms. Joint semialgebraic certificates handle dominant vectors that avoid zero on the compact toric closure. Binomial-stratified certificates handle dominant zero loci contained in finitely many character strata, because exact character hits are computable and Baker separation gives polynomial distance from the strata away from those hits.

A scalar expression such as
\[
X+Y-1
\]
on a genuinely rank-two orbit is different. Its zero locus is not a finite union of binomial strata. Along
\[
(X,Y)=(\alpha^n,\beta^n)
\]
one would need effective lower bounds for
\[
|\alpha^n+\beta^n-1|
\]
outside exact zeros. Such non-binomial toric approximation is not supplied by the binomial-stratified layer.

Thus the compiler is sound and incomplete: success produces a verified exact zero set, while failure means that the available certificate mechanisms did not certify the instance.

\section{Conclusion}

We have developed a certificate framework for effective zero-free tails in scalar and vector linear recurrence sequences. The main novelty is the vector toric viewpoint: for Simultaneous Skolem and Subspace Orbit, it is enough to separate the dominant vector from zero. This can succeed even when scalar coordinatewise methods fail.

The framework includes:
\begin{enumerate}
\item global vector dominant-shell normal form;
\item torsion saturation and toric coset closures;
\item certified-tail soundness;
\item joint semialgebraic and SOS certificates;
\item Bezout--toric certificates;
\item scalar Baker character-polynomial atoms;
\item quotient-character reduction;
\item binomial-stratified lower bounds;
\item Baker separation from exact binomial hits;
\item automatic binomial extraction from dominant elimination ideals;
\item finite certificate trees;
\item a partial Factor--Reduce--Certify compiler with a recursive binomial-stratified mode;
\item applications to Simultaneous Skolem, Subspace Orbit, linear-loop reachability, algebraic target hitting, and ultimate positivity;
\item a high-rank vector family showing genuine advantage over scalar methods.
\end{enumerate}

The resulting framework gives a rigorous, checkable, and partially automatic mechanism for proving effective Skolem tails in high-rank, high-order families. Its strongest automatic layer is the binomial-stratified class: semialgebraic Lojasiewicz inequalities control vanishing near the strata, Baker separation controls the actual toric orbit away from exact hits, and recursive restriction handles the exact-hit progressions.

\begin{thebibliography}{99}

\bibitem{Baker}
A. Baker.
\emph{Transcendental Number Theory}.
Cambridge University Press, 1975.

\bibitem{BakerWuestholz}
A. Baker and G. Wuestholz.
\emph{Logarithmic Forms and Diophantine Geometry}.
Cambridge University Press, 2007.

\bibitem{BacikOrder4}
P. Bacik.
Completing the picture for the Skolem Problem on order-4 linear recurrence sequences.
arXiv:2409.01221.

\bibitem{BacikPadic}
P. Bacik, J. Ouaknine, D. Purser, and J. Worrell.
On the \(p\)-adic Skolem Problem.
arXiv:2504.14413.

\bibitem{BacikVaronka}
P. Bacik and A. Varonka.
On the Subspace Orbit Problem and the Simultaneous Skolem Problem.
arXiv:2601.18349.

\bibitem{Basu}
S. Basu.
Algorithms in real algebraic geometry: a survey.
arXiv:1409.1534.

\bibitem{Bilu}
Y. Bilu.
Skolem Problem for linear recurrence sequences with 4 dominant roots, after Mignotte, Shorey, Tijdeman, Vereshchagin and Bacik.
arXiv:2501.16290.

\bibitem{Schanuel}
Y. Bilu, F. Luca, J. Nieuwveld, J. Ouaknine, D. Purser, and J. Worrell.
Skolem meets Schanuel.
arXiv:2204.13417.

\bibitem{Chonev}
V. Chonev, J. Ouaknine, and J. Worrell.
The Orbit Problem in higher dimensions.
In \emph{Proceedings of LICS 2013}.

\bibitem{ESS}
J.-H. Evertse, H. P. Schlickewei, and W. M. Schmidt.
Linear equations in variables which lie in a multiplicative group.
\emph{Annals of Mathematics} 155 (2002), 807--836.

\bibitem{EisenbudSturmfels}
D. Eisenbud and B. Sturmfels.
Binomial ideals.
\emph{Duke Mathematical Journal} 84 (1996), 1--45.

\bibitem{Ge}
G. Ge.
\emph{Algorithms Related to Multiplicative Representations of Algebraic Numbers}.
PhD thesis, University of California, Berkeley, 1993.

\bibitem{Jeronimo}
G. Jeronimo, D. Perrucci, and E. Tsigaridas.
On the minimum of a polynomial function on a basic closed semialgebraic set and applications.
arXiv:1112.0544.

\bibitem{KannanLipton}
R. Kannan and R. Lipton.
The Orbit Problem is decidable.
In \emph{Proceedings of STOC 1980}.

\bibitem{Mignotte}
M. Mignotte, T. N. Shorey, and R. Tijdeman.
The distance between terms of an algebraic recurrence sequence.
\emph{Journal fuer die reine und angewandte Mathematik}.

\bibitem{SchlickeweiSchmidt}
H. P. Schlickewei and W. M. Schmidt.
The number of solutions of polynomial-exponential equations.
\emph{Compositio Mathematica} 120 (2000), 193--225.

\bibitem{ShoreyTijdeman}
T. N. Shorey and R. Tijdeman.
\emph{Exponential Diophantine Equations}.
Cambridge University Press, 1986.

\bibitem{Vereshchagin}
N. K. Vereshchagin.
The problem of the occurrence of zero in a linear recursive sequence.
\emph{Mathematical Notes}.

\end{thebibliography}

\end{document}
